% Auto-generated by scripts/analyze_github_project_data.mjs
\subsection{GitHub Project Corpus and Sampling Funnel}
\label{subsec:github-project-corpus-sampling}
The GitHub corpus is organized as a funnel rather than a flat awesome list. The raw discovery layer contains 416 timestamp-indexed GitHub captures. The classification layer contains 416 repositories with category, theme, stack, and time-slice annotations. The paper-facing analysis layer currently contains 129 repositories that receive public project reports and model-card pages. Within the classified corpus, 71 repositories are strict evolution-theme repositories, while 166 match a broader evolution-related keyword set including self-improvement, reflection, recursive improvement, prompt evolution, and code evolution.

\begin{table}[t]
\centering
\caption{GitHub project corpus funnel.}
\label{tab:github-project-corpus-funnel}
\begin{tabular}{lrp{0.55\linewidth}}
\toprule
Layer & Count & Meaning \\
\midrule
Raw captures & 416 & Timestamp-indexed records from \texttt{raw-github/}. \\
Classified repositories & 416 & Rows with category, theme, stack, evidence, and time-slice fields. \\
Analyzed model-card projects & 129 & Repositories promoted into site and paper-facing project analysis. \\
Strict evolution theme & 71 & Classified rows whose base theme is \texttt{evolution}. \\
Broad evolution-related & 166 & Rows matching evolution, self-improvement, reflection, or recursive-improvement signals. \\
\bottomrule
\end{tabular}
\end{table}

\begin{table}[t]
\centering
\caption{Raw GitHub category and theme distribution.}
\label{tab:github-category-theme-distribution}
\begin{tabular}{lrlr}
\toprule
Category & Count & Theme & Count \\
\midrule
框架/framework & 118 & evaluation & 89 \\
评测/evaluation & 96 & memory & 77 \\
教程/tutorial & 75 & evolution & 71 \\
工具/tool & 56 & framework & 43 \\
应用/application & 47 & education-list & 35 \\
论文代码/paper-code & 24 & research-agent & 31 \\
 &  & prompt-optimization & 26 \\
 &  & skill & 20 \\
 &  & coding-agent & 17 \\
 &  & workflow-automation & 6 \\
 &  & safety & 1 \\
\bottomrule
\end{tabular}
\end{table}

The time analysis deliberately separates repository creation from activity timestamps. The raw corpus time-slice field is an activity/content timestamp extracted from GitHub captures, so it measures when the captured page exposed recent activity. The analyzed-project timeline uses GitHub API \texttt{created\_at} where available. When GitHub API creation time is unavailable, the table marks a weaker \texttt{local\_git\_first\_commit} fallback; this is a first-observed mirror timestamp, not a release-date claim. In the current run, 49 of 129 analyzed projects have either verified creation time or a local fallback.

\begin{table*}[t]
\centering
\small
\caption{First 24 analyzed projects by verified creation month or local first-observed fallback.}
\label{tab:analyzed-project-release-timeline}
\begin{tabular}{lllp{0.21\linewidth}p{0.30\linewidth}}
\toprule
Month & Source & Repository & Category & Evolution pattern \\
\midrule
2022-09 & github\_api & \texttt{carperai/openelm} & 进化式 Prompt 优化 & 进化/搜索循环 → 评估器/打分器 \\
2023-01 & github\_api & \texttt{stanfordnlp/dspy} & 声明式 Prompt 优化 & 反馈-精炼 → 进化/搜索循环 → 评估器/打分器 \\
2023-03 & github\_api & \texttt{Significant-Gravitas/AutoGPT} & 自主 Agent 平台 & 智能体编排 \\
2023-03 & github\_api & \texttt{camel-ai/camel} & 角色扮演 Agent 框架 & 智能体编排 → 反馈-精炼 \\
2023-03 & github\_api & \texttt{noahshinn/reflexion} & 反思记忆 & 进化/搜索循环 → 反思记忆 → 反馈-精炼 → 评估器/打分器 → 训练/数据循环 \\
2023-03 & github\_api & \texttt{madaan/self-refine} & 反馈精炼 & 反馈-精炼 \\
2023-04 & local\_git\_first\_commit & \texttt{automl/auto-sklearn} & AutoML 框架 & 进化/搜索循环 → 评估器/打分器 \\
2023-04 & github\_api & \texttt{microsoft/CoML} & ML 知识库驱动 & 反馈-精炼 → 评估器/打分器 \\
2023-06 & github\_api & \texttt{FoundationAgents/MetaGPT} & 多 Agent 协作框架 & 智能体编排 → 反馈-精炼 \\
2023-07 & github\_api & \texttt{aiwaves-cn/agents} & 数据驱动 Agent 进化 & 进化/搜索循环 → 评估器/打分器 → 智能体编排 \\
2023-08 & github\_api & \texttt{langchain-ai/langgraph} & 图式 Agent 编排 & 智能体编排 \\
2023-08 & github\_api & \texttt{microsoft/autogen} & 多 Agent 对话框架 & 智能体编排 \\
2023-10 & github\_api & \texttt{google-deepmind/opro} & LLM 作为优化器 & 进化/搜索循环 → 评估器/打分器 \\
2023-10 & github\_api & \texttt{crewAIInc/crewAI} & 多 Agent 协作框架 & 智能体编排 \\
2023-11 & github\_api & \texttt{google-deepmind/funsearch} & 进化式数学发现 & 进化/搜索循环 → 评估器/打分器 \\
2024-03 & github\_api & \texttt{All-Hands-AI/OpenHands} & AI 软件开发平台 & 智能体编排 \\
2024-04 & github\_api & \texttt{princeton-nlp/SWE-agent} & 软件工程 Agent & 反馈-精炼 → 评估器/打分器 \\
2024-07 & github\_api & \texttt{shengranhu/adas} & Agent 架构自动搜索 & 进化/搜索循环 → 智能体编排 → 评估器/打分器 \\
2024-08 & local\_git\_first\_commit & \texttt{alfa-group/tutorial\_gp\_llm} & GP+LLM 教学 & 进化/搜索循环 → 评估器/打分器 \\
2024-09 & local\_git\_first\_commit & \texttt{OpenBMB/AgentVerse} & 多 Agent 仿真平台 & 智能体编排 → 反思记忆 \\
2024-10 & local\_git\_first\_commit & \texttt{siyuyuan/evoagent} & 进化式多 Agent 系统 & 进化/搜索循环 → 智能体编排 \\
2024-11 & local\_git\_first\_commit & \texttt{xiaofangxd/LLM\_EA} & LLM+EA 交叉综述 & 文献综述 \\
2025-03 & local\_git\_first\_commit & \texttt{wuxingyu-ai/LLM4EC} & LLM+EC 交叉综述 & 文献综述 \\
2025-05 & github\_api & \texttt{DeepAuto-AI/automl-agent} & 多 Agent AutoML & 智能体编排 → 进化/搜索循环 → 评估器/打分器 \\
\bottomrule
\end{tabular}
\end{table*}

\begin{table}[t]
\centering
\caption{Top strict evolution-theme repositories by local star snapshot.}
\label{tab:top-evolution-theme-repos}
\begin{tabular}{lrrp{0.42\linewidth}}
\toprule
Repository & Stars & Time & Description \\
\midrule
\texttt{aiming-lab/AutoResearchClaw} & 12600 & 2026-05 & AutoResearchClaw 是从 research idea 到 paper 的自主/协作式科研 agent 管线，结合多 agent debate、实验沙箱、claim verification、HITL co-pilot、MetaClaw cross-run learning 和 ARC-Bench。 \\
\texttt{aden-hive/hive} & 10400 & 2026-05 & Aden Hive 是生产 AI 的 multi-agent harness，强调状态持久化、崩溃恢复、成本控制、审计轨迹、MCP 工具集成和失败驱动的 graph evolution。 \\
\texttt{metauto-ai/gptswarm} & 998 & 2026-05 &  The First Self-Improving agents with RL / Prompting Optimization \\
\texttt{adolfousier/opencrabs} & 755 & 2026-05 & OpenCrabs 是受 OpenClaw 启发的单二进制多渠道 AI agent，强调本地 brain files、memory search、custom commands、cron/background jobs 和 self-update 组成的自改进循环。 \\
\texttt{human-agent-society/coral} & 667 & 2026-05 & CORAL 是面向 open-ended discovery 的多代理自主演化基础设施，用隔离 git worktrees、共享状态目录、grader daemon、heartbeat prompt 和多 runtime 集成推动 agent teams 连续改进。 \\
\texttt{a-evo-lab/a-evolve} & 552 & 2026-05 & A-Evolve 是通用 self-improving agent 基础设施：给定 base agent、benchmark 和 evolution algorithm，就把 prompt、skills、memory 等 agent workspace 文件作为可变状态进行迭代。 \\
\texttt{thudm/webrl} & 524 & 2026-05 & Building Open LLM Web Agents with Self-Evolving Online Curriculum RL \\
\texttt{ecnu-icalk/autoskill} & 424 & 2026-05 & AutoSkill: Experience-Driven Lifelong Learning via Skill Self-Evolution \\
\texttt{china-qijizhifeng/agentic-Harness-engineering} & 391 & 2026-05 & Agentic Harness Engineering treats prompts, tools, middleware, skills, sub-agents, memory, and evaluator scaffolds as evolvable harness surfaces while the base model stays fixed. \\
\texttt{feiliu36/eoh} & 319 & unknown & Evolution of Heuristics \\
\texttt{meituan/EvoCUA} & 317 & 2026-05 & 美团 EvoCUA 计算机使用 Agent 项目，发布 EvoCUA-32B/8B 并在 OSWorld、WindowsAgentArena 等 GUI 自动化评测上报告结果。 \\
\texttt{brain-research/guided-evolutionary-strategies} & 273 & unknown & Guided Evolutionary Strategies \\
\bottomrule
\end{tabular}
\end{table}
